% 蓝、红、绿标题分别搭配同色系浅底；色带略深，正文区域保留清楚的底色。
\definecolor{flowInk}{HTML}{08264A}
\definecolor{flowBlueHead}{HTML}{D1E3F1}
\definecolor{flowBlueBody}{HTML}{E5EFF7}
\definecolor{flowBlueEdge}{HTML}{8EB9D5}
\definecolor{flowBlueTitle}{HTML}{14538B}
\definecolor{flowWarmHead}{HTML}{EDCDC8}
\definecolor{flowWarmBody}{HTML}{F6E5E2}
\definecolor{flowWarmEdge}{HTML}{BB827D}
\definecolor{flowWarmTitle}{HTML}{93332D}
\definecolor{flowGreenHead}{HTML}{CEE4DA}
\definecolor{flowGreenBody}{HTML}{E5F0E9}
\definecolor{flowGreenEdge}{HTML}{78A591}
\definecolor{flowGreenTitle}{HTML}{267A63}
\definecolor{flowNote}{HTML}{647889}
\tikzset{
  paper flow/.style={>=Stealth,line cap=butt,line join=miter,
    text=flowInk,font=\linespread{1}\fontsize{9.5}{12}\selectfont},
  flow arrow/.style={->,draw=flowBlueTitle,line width=.8pt},
  flow label/.style={fill=none,inner sep=2pt,align=center,
    font=\linespread{1}\fontsize{9.5}{12}\selectfont}
}

% 参数：名称、中心、宽(cm)、高(cm)、配色、标题、说明。
% 标题色带与正文分开排版；所有框均为直角，说明文字无白色遮罩。
\newcommand{\PaperFlowCard}[7]{%
  \node[rectangle,inner sep=0pt,minimum width=#3cm,minimum height=#4cm,
    fill=flow#5Body,draw=none] (#1) at #2 {};
  \path[fill=flow#5Head] (#1.north west)
    rectangle ($(#1.north east)+(0,-.49)$);
  \draw[flow#5Edge,line width=.65pt] (#1.north west) rectangle (#1.south east);
  \node[anchor=north,inner sep=0pt,text=flow#5Title,
    font=\linespread{1}\fontsize{10.5}{12.5}\selectfont\bfseries]
    at ($(#1.north)+(0,-.075)$) {#6};
  \node[anchor=north,inner sep=0pt,align=center,
    text width=\dimexpr#3cm-4mm\relax,font=\linespread{1}\fontsize{9.5}{12}\selectfont]
    at ($(#1.north)+(0,-.66)$) {#7};
}

% 第四问保留已确认的分支顺序，使用相同的浅色与直角框。
\colorlet{flowLine}{flowBlueTitle}
\colorlet{flowAccent}{flowBlueTitle}
\tikzset{
  flow process/.style={rectangle,draw=flowBlueEdge,fill=flowBlueBody,line width=.65pt,
    inner xsep=2mm,inner ysep=2mm,align=center,
    font=\linespread{1}\fontsize{10}{12.5}\selectfont}
}
